% Generated draft: requires mathematical and source review.
\documentclass[11pt]{article}
\usepackage[margin=25mm]{geometry}
\usepackage{fontspec}
\setmainfont{Latin Modern Roman}
\usepackage{amsmath,amssymb,mathtools,unicode-math}
\setmathfont{Latin Modern Math}
\usepackage{enumitem,needspace,xurl,hyperref}
\setlength{\parindent}{0pt}
\setlength{\parskip}{5pt}
\setlength{\emergencystretch}{3em}
\newcommand{\N}{\mathbb{N}}
\newcommand{\Z}{\mathbb{Z}}
\newcommand{\Q}{\mathbb{Q}}
\newcommand{\R}{\mathbb{R}}
\newcommand{\C}{\mathbb{C}}
\newcommand{\F}{\mathbb{F}}
\newcommand{\A}{\mathbb{A}}
\newcommand{\PP}{\mathbb{P}}
\newcommand{\E}{\mathbb{E}}
\newcommand{\eps}{\varepsilon}
\newcommand{\dd}{\,\mathrm{d}}
\begin{document}
\section*{UnsolvedMath ID: 30002525}
\textbf{Source problem code:} catalogue entry
\subsection{Definitions and mathematical statement}
A complex projective variety is a complex algebraic variety admitting a closed embedding into some complex projective space $\PP^n_{\C}$. It is irreducible if it is not the union of two proper Zariski-closed subsets. In this draft, $\pi_1(X(\C),x)$ denotes the topological fundamental group of the complex points with their usual analytic topology, based at $x\in X(\C)$.

The supplied question is: ``Which groups occur as fundamental groups of complex projective varieties when the varieties are required to be irreducible and their singularities are suitably controlled?''

To display the unresolved qualification explicitly, let $\mathcal P(X)$ stand for the intended condition on the singularities of $X$. For a specified such condition, the classification target would be
\[
\mathcal G_{\mathcal P}
=
\left\{[G]_{\cong}\;\middle|\;
\begin{gathered}
G\text{ is a group, and there exist an irreducible complex}\\
\text{projective variety }X\text{ and }x\in X(\C)\text{ such that}\\
\mathcal P(X)\quad\text{and}\quad G\cong\pi_1(X(\C),x)
\end{gathered}
\right\},
\]
where $[G]_{\cong}$ denotes the isomorphism class of $G$. The problem asks for a characterization of the groups in this collection.

The source does not specify $\mathcal P$: in particular, it gives no permitted local singularity types or other precise singularity restrictions. Consequently, this is only a conditional formulation of the classification question; it does not determine a single classification problem until the intended singularity condition is recovered. No dimension restriction is supplied.
\subsection{Short English statement}
Determine which groups can be realized as fundamental groups of irreducible complex projective varieties subject to the intended restrictions on their singularities. Those restrictions are unspecified in the supplied record, so a precise classification target remains unresolved. This draft uses the usual topological fundamental group.
\subsection{Sources}
\begin{itemize}
\item Supplied local UnsolvedMath registry record: permanent ID $30002525$, title \emph{Fundamental Groups of Singular Irreducible Projective Varieties}, source problem code \texttt{catalogue entry}, in the Oberwolfach Reports Open Problems collection. The proposed year is recorded as 2014; no individual proposer is supplied.
\item Supplied report attribution: \emph{Mini-Workshop: Kähler Groups} (2014), Oberwolfach Reports, DOI 10.4171/owr/2014/09, pp.\ 531--553. \url{https://doi.org/10.4171/owr/2014/09}. The supplied background describes these reports as documenting mathematical workshops and their open-problem discussions.
\item Supplied literature source: Kapovich, 2013. \url{https://doi.org/10.1007/s00222-013-0453-4}. No title or more detailed bibliographic information is supplied for this reference.
\item Supplied literature source: Oberwolfach Report 9, 2014. \url{https://doi.org/10.4171/owr/2014/09}.
\item Literature assessment dated exactly \texttt{2026-08-21}: the supplied curation labels the problem open and states that projective realizations of finitely presented groups with singularities are known, while the requested simultaneous control of singularities and irreducibility is unresolved. This is supplied dated curation derived from the OWR material, not an independently verified literature assessment. The record expressly calls for independent verification of time-sensitive status and cited claims.
\end{itemize}
\end{document}
